% =============================================================
%  03_suggerimenti.tex — Suggerimenti per esercizi ★★★ e superiori
% =============================================================
\chapter*{Suggerimenti}
\addcontentsline{toc}{chapter}{Suggerimenti}
\markboth{Suggerimenti}{Suggerimenti}
\setcounter{hintsec}{0}  % reset contatore sezioni hint

\noindent\small\color{black!65}
Suggerimenti per gli esercizi \treStelle\ o superiore.
Prima di leggere, fai almeno un tentativo autonomo: anche una soluzione parziale è più utile della risposta immediata.
\color{black}

\vspace{8pt}

% =============================================================
\hintsection{Fondamenti}
% =============================================================

\label{hint:17}
\begin{hintbox}{17}{Somma delle cifre}
\small\color{black!55}\hyperref[ex:17]{$\leftarrow$ Esercizio}\normalsize\color{black}

Estrai le cifre una alla volta con \fn{mod} e \fn{floor}:
\fn{(mod n 10)} da' l'ultima cifra; \fn{(floor n 10)} rimuove
l'ultima cifra. Continua finche' \fn{n} e' zero.

\small\color{black!55}\hyperref[sol:17]{$\rightarrow$ Soluzione}\normalsize\color{black}
\end{hintbox}

% =============================================================
\hintsection{Liste}
% =============================================================

\label{hint:40}
\begin{hintbox}{40}{Append}
\small\color{black!55}\hyperref[ex:40]{$\leftarrow$ Esercizio}\normalsize\color{black}

Casi: se \fn{lst1} e' vuota, restituisci \fn{lst2}.
Altrimenti costruisci con \fn{cons}: il nuovo primo elemento
e' \fn{(car lst1)}, il resto e' la ricorsione su \fn{(cdr lst1)}
e \fn{lst2}.
Non modificare \fn{lst2}: non servono copie, \fn{cons} crea
una nuova struttura automaticamente.

\small\color{black!55}\hyperref[sol:40]{$\rightarrow$ Soluzione}\normalsize\color{black}
\end{hintbox}

\label{hint:41}
\begin{hintbox}{41}{Remove-first}
\small\color{black!55}\hyperref[ex:41]{$\leftarrow$ Esercizio}\normalsize\color{black}

Caso base 1: lista vuota $\to$ \fn{NIL}.
Caso base 2: \fn{(car lst)} uguale a \fn{item} $\to$ restituisci
\fn{(cdr lst)} (salta solo questo elemento).
Caso ricorsivo: \fn{cons} di \fn{(car lst)} con la ricorsione
sul resto. Non rimuovere altre occorrenze: basta non continuare
a cercare dopo aver trovato la prima.

\small\color{black!55}\hyperref[sol:41]{$\rightarrow$ Soluzione}\normalsize\color{black}
\end{hintbox}

\label{hint:49}
\begin{hintbox}{49}{Run-length encoding}
\small\color{black!55}\hyperref[ex:49]{$\leftarrow$ Esercizio}\normalsize\color{black}

\hintschema{Encode:} tieni traccia dell'elemento corrente e del contatore.
Quando l'elemento cambia (o la lista e' finita), aggiungi
la coppia \fn{(contatore elemento)} al risultato.
Una struttura con accumulatori e' piu' naturale di una
puramente ricorsiva.

\hintschema{Decode:} per ogni coppia \fn{(n elem)}, genera \fn{n}
copie di \fn{elem}. Usa un loop interno o una funzione ausiliaria
che costruisce \fn{n} ripetizioni.

\small\color{black!55}\hyperref[sol:49]{$\rightarrow$ Soluzione}\normalsize\color{black}
\end{hintbox}

\label{hint:50}
\begin{hintbox}{50}{Remove-duplicates*}
\small\color{black!55}\hyperref[ex:50]{$\leftarrow$ Esercizio}\normalsize\color{black}

Visita la lista da sinistra. Per ogni elemento, controllane
la presenza nel risultato accumulato (o nella parte gia'
processata). Se non e' presente, includilo; altrimenti saltalo.
Usa \fn{member} (o la tua \fn{member*}) per il controllo.

\small\color{black!55}\hyperref[sol:50]{$\rightarrow$ Soluzione}\normalsize\color{black}
\end{hintbox}

% =============================================================
\hintsection{Ricorsione avanzata}
% =============================================================

\label{hint:52}
\begin{hintbox}{52}{Reverse tail-recursive}
\small\color{black!55}\hyperref[ex:52]{$\leftarrow$ Esercizio}\normalsize\color{black}

Definire con \fn{labels} una funzione ausiliaria con due parametri:
\fn{lst} (la lista da rovesciare) e \fn{acc} (l'accumulatore, inizialmente \fn{NIL}).
Ad ogni passo: se \fn{lst} è vuota restituire \fn{acc};
altrimenti richiamare la funzione con \fn{(cdr lst)} e
\fn{(cons (car lst) acc)}.
La chiamata ricorsiva è l'ultima operazione: nessun calcolo al ritorno.

\small\color{black!55}\hyperref[sol:52]{$\rightarrow$ Soluzione}\normalsize\color{black}
\end{hintbox}

\label{hint:54}
\begin{hintbox}{54}{Fibonacci O(n)}
\small\color{black!55}\hyperref[ex:54]{$\leftarrow$ Esercizio}\normalsize\color{black}

Usa due accumulatori \fn{a} e \fn{b} che rappresentano
$F(i)$ e $F(i+1)$. A ogni passo: il nuovo \fn{a} diventa
il vecchio \fn{b}; il nuovo \fn{b} diventa \fn{a+b}.
Alla fine, \fn{a} e' $F(n)$.

Caso base: se \fn{n = 0}, restituisci \fn{a} (inizializzato a 0).
Inizia con \fn{(iter n 0 1)}.

\small\color{black!55}\hyperref[sol:54]{$\rightarrow$ Soluzione}\normalsize\color{black}
\end{hintbox}

\label{hint:55}
\begin{hintbox}{55}{Flatten*}
\small\color{black!55}\hyperref[ex:55]{$\leftarrow$ Esercizio}\normalsize\color{black}

Tratta ogni elemento della lista come il nodo di un albero binario
implicito: se e' un \textit{atomo}, e' una foglia (mettila nel risultato);
se e' una \textit{lista}, e' un sottoalbero (appiattiscila ricorsivamente).

Schema:
\begin{lstlisting}
(defun flatten* (tree)
  (cond
    ((null tree) '())
    ((atom tree) (list tree))
    (t (append (flatten* (car tree))
               (flatten* (cdr tree))))))
\end{lstlisting}
\hintschema{Attenzione:} ricorrere su \fn{car} e \fn{cdr} separatamente,
non su \fn{(car tree)} e \fn{(cdr tree)} come lista.

\small\color{black!55}\hyperref[sol:55]{$\rightarrow$ Soluzione}\normalsize\color{black}
\end{hintbox}

\label{hint:56}
\begin{hintbox}{56}{Tree-depth}
\small\color{black!55}\hyperref[ex:56]{$\leftarrow$ Esercizio}\normalsize\color{black}

La profondita' di una lista vuota e' 0.
La profondita' di una lista non vuota e' il massimo tra:
\begin{itemize}[leftmargin=2em]
  \item la profondita' del \fn{car} $+1$ (se il \fn{car} e' una lista);
  \item la profondita' del \fn{cdr}.
\end{itemize}
Usa \fn{max} per confrontare i due rami.

\small\color{black!55}\hyperref[sol:56]{$\rightarrow$ Soluzione}\normalsize\color{black}
\end{hintbox}

\label{hint:60}
\begin{hintbox}{60}{Sostituzione profonda}
\small\color{black!55}\hyperref[ex:60]{$\leftarrow$ Esercizio}\normalsize\color{black}

Struttura simile a \fn{flatten*}: ricorri su \fn{car} e \fn{cdr}.
Se il \fn{car} e' uguale a \fn{old}, sostituiscilo con \fn{new};
se e' una lista, ricorri dentro; altrimenti lascialo invariato.

\small\color{black!55}\hyperref[sol:60]{$\rightarrow$ Soluzione}\normalsize\color{black}
\end{hintbox}

\label{hint:61}
\begin{hintbox}{61}{Power-fast}
\small\color{black!55}\hyperref[ex:61]{$\leftarrow$ Esercizio}\normalsize\color{black}

Tre casi: \fn{exp = 0} $\to$ 1; \fn{exp} pari $\to$ calcola
\fn{(power-fast base (/ exp 2))} e metti al quadrato;
\fn{exp} dispari $\to$ moltiplica per \fn{base} e richiama
con \fn{(1- exp)}.

Con \fn{let} puoi evitare la doppia chiamata ricorsiva
nel caso pari:
\begin{lstlisting}
(let ((half (power-fast base (/ exp 2))))
  (* half half))
\end{lstlisting}

\small\color{black!55}\hyperref[sol:61]{$\rightarrow$ Soluzione}\normalsize\color{black}
\end{hintbox}

\label{hint:64}
\begin{hintbox}{64}{Crivello di Eratostene}
\small\color{black!55}\hyperref[ex:64]{$\leftarrow$ Esercizio}\normalsize\color{black}

\begin{enumerate}[leftmargin=2em]
  \item Crea un array booleano \fn{sieve} di dimensione $n+1$,
        inizializzato a \fn{T} (tutti candidati).
  \item Marca 0 e 1 come non-primi.
  \item Per ogni \fn{i} da 2 a $\sqrt{n}$: se \fn{(aref sieve i)}
        e' \fn{T}, marca tutti i multipli di \fn{i} (da $i^2$)
        come \fn{NIL}.
  \item Raccogli tutti gli indici per cui il valore e' ancora \fn{T}.
\end{enumerate}

\small\color{black!55}\hyperref[sol:64]{$\rightarrow$ Soluzione}\normalsize\color{black}
\end{hintbox}

\label{hint:65}
\begin{hintbox}{65}{Fattori primi}
\small\color{black!55}\hyperref[ex:65]{$\leftarrow$ Esercizio}\normalsize\color{black}

Prova a dividere \fn{n} per ogni intero \fn{d} partendo da 2.
Se \fn{(mod n d) = 0}, \fn{d} e' un fattore: aggiungilo al risultato
e continua con \fn{n/d} (non avanzare \fn{d}!).
Se \fn{d} non divide, avanza \fn{d} di 1.
Termina quando \fn{d * d > n}: se \fn{n > 1} e' rimasto
un fattore primo.

\small\color{black!55}\hyperref[sol:65]{$\rightarrow$ Soluzione}\normalsize\color{black}
\end{hintbox}

\label{hint:67}
\begin{hintbox}{67}{Raggruppa in blocchi}
\small\color{black!55}\hyperref[ex:67]{$\leftarrow$ Esercizio}\normalsize\color{black}

Scrivi una funzione ausiliaria \fn{prendi-n lst n} che restituisce
\fn{(primi-n resto)}: i primi \fn{n} elementi e il resto della lista.
Poi \fn{raggruppa} chiama \fn{prendi-n}, prende il primo gruppo,
e ricorre sul resto.

\small\color{black!55}\hyperref[sol:67]{$\rightarrow$ Soluzione}\normalsize\color{black}
\end{hintbox}

% =============================================================
\hintsection{Higher-order e closure}
% =============================================================

\label{hint:75}
\begin{hintbox}{75}{Cifre in intero}
\small\color{black!55}\hyperref[ex:75]{$\leftarrow$ Esercizio}\normalsize\color{black}

Usa \fn{reduce} con una funzione che, dato l'accumulatore \fn{acc}
e la cifra corrente \fn{d}, calcola \fn{(+ (* acc 10) d)}.
Con \fn{:initial-value 0} funziona anche per lista vuota.

\small\color{black!55}\hyperref[sol:75]{$\rightarrow$ Soluzione}\normalsize\color{black}
\end{hintbox}

\label{hint:77}
\begin{hintbox}{77}{Regola di Horner}
\small\color{black!55}\hyperref[ex:77]{$\leftarrow$ Esercizio}\normalsize\color{black}

La regola di Horner valuta $p(x) = c_n x^n + \ldots + c_0$
come $(\ldots((c_n \cdot x + c_{n-1}) \cdot x + c_{n-2})\ldots + c_0)$.

Con \fn{reduce}: la funzione riceve l'accumulatore \fn{acc}
(che parte da $c_n$) e il coefficiente corrente \fn{c},
e calcola \fn{(+ (* acc x) c)}.
I coefficienti devono essere nell'ordine dal grado piu' alto
al piu' basso.

\small\color{black!55}\hyperref[sol:77]{$\rightarrow$ Soluzione}\normalsize\color{black}
\end{hintbox}

\label{hint:83}
\begin{hintbox}{83}{Raggruppa per chiave}
\small\color{black!55}\hyperref[ex:83]{$\leftarrow$ Esercizio}\normalsize\color{black}

Per ogni elemento, calcola la chiave. Cerca nella alist risultato
se quella chiave esiste gia'. Se si', aggiungici l'elemento;
se no, crea una nuova coppia \fn{(chiave elemento)}.
Usa \fn{assoc} per la ricerca e aggiorna la alist con \fn{setf}/\fn{aref}
o ricostruendola.

\small\color{black!55}\hyperref[sol:83]{$\rightarrow$ Soluzione}\normalsize\color{black}
\end{hintbox}

\label{hint:85}
\begin{hintbox}{85}{Compose}
\small\color{black!55}\hyperref[ex:85]{$\leftarrow$ Esercizio}\normalsize\color{black}

Una closure cattura \fn{f} e \fn{g} nell'ambiente lessicale.
Restituisce una \fn{lambda} che prende \fn{x}, applica prima
\fn{g}, poi \fn{f}:
\begin{lstlisting}
(lambda (x) (funcall f (funcall g x)))
\end{lstlisting}
Ricorda che le funzioni si passano con \fn{\#'}.

\small\color{black!55}\hyperref[sol:85]{$\rightarrow$ Soluzione}\normalsize\color{black}
\end{hintbox}

\label{hint:86}
\begin{hintbox}{86}{Partial application}
\small\color{black!55}\hyperref[ex:86]{$\leftarrow$ Esercizio}\normalsize\color{black}

Cattura \fn{funzione} e \fn{argomenti} nella closure.
La lambda restituita prende \fn{\&rest rest-args} e chiama
\fn{funzione} con la lista \fn{(append argomenti rest-args)}.
Usa \fn{apply} per passare la lista combinata:
\begin{lstlisting}
(lambda (&rest rest-args)
  (apply funzione (append argomenti rest-args)))
\end{lstlisting}

\small\color{black!55}\hyperref[sol:86]{$\rightarrow$ Soluzione}\normalsize\color{black}
\end{hintbox}

\label{hint:87}
\begin{hintbox}{87}{Closure con stato}
\small\color{black!55}\hyperref[ex:87]{$\leftarrow$ Esercizio}\normalsize\color{black}

Lo stato vive in un \fn{let} attorno alla \fn{lambda}.
Ogni chiamata alla lambda legge e aggiorna la variabile.
\begin{lstlisting}
(let ((curr start))
  (lambda ()
    (prog1 curr
      (incf curr step))))
\end{lstlisting}
\fn{prog1} restituisce il valore prima dell'incremento.

\small\color{black!55}\hyperref[sol:87]{$\rightarrow$ Soluzione}\normalsize\color{black}
\end{hintbox}

\label{hint:89}
\begin{hintbox}{89}{Generatore Fibonacci con closure}
\small\color{black!55}\hyperref[ex:89]{$\leftarrow$ Esercizio}\normalsize\color{black}

Due variabili nello stesso \fn{let} condividono la stessa closure:
\begin{lstlisting}
(let ((a 0) (b 1))
  (lambda ()
    (prog1 a
      (psetf a b b (+ a b)))))
\end{lstlisting}
\fn{psetf} aggiorna le variabili simultaneamente (senza
che il nuovo valore di \fn{a} influenzi il calcolo di \fn{b}).

\small\color{black!55}\hyperref[sol:89]{$\rightarrow$ Soluzione}\normalsize\color{black}
\end{hintbox}

\label{hint:91}
\begin{hintbox}{91}{Memoize}
\small\color{black!55}\hyperref[ex:91]{$\leftarrow$ Esercizio}\normalsize\color{black}

La struttura base:
\begin{lstlisting}
(let ((cache (make-hash-table :test #'equal)))
  (lambda (&rest args)
    (multiple-value-bind (val found)
        (gethash args cache)
      (if found
          val
          (setf (gethash args cache)
                (apply funzione args))))))
\end{lstlisting}
La chiave della cache e' la lista degli argomenti: usa
\fn{:test \#'equal} per confrontare liste.
\fn{multiple-value-bind} permette di distinguere ``valore NIL
trovato'' da ``chiave non presente''.

\small\color{black!55}\hyperref[sol:91]{$\rightarrow$ Soluzione}\normalsize\color{black}
\end{hintbox}

% =============================================================
\hintsection{Strutture dati}
% =============================================================

\label{hint:94}
\begin{hintbox}{94}{Frequenze con alist}
\small\color{black!55}\hyperref[ex:94]{$\leftarrow$ Esercizio}\normalsize\color{black}

Per ogni elemento: cerca nella alist con \fn{assoc}.
Se la coppia esiste, incrementa il \fn{cdr}; altrimenti,
aggiungi una nuova coppia \fn{(elemento . 1)}.
Con \fn{setf} puoi modificare il \fn{cdr} di una coppia
trovata con \fn{assoc}: \fn{(setf (cdr coppia) (1+ (cdr coppia)))}.

\small\color{black!55}\hyperref[sol:94]{$\rightarrow$ Soluzione}\normalsize\color{black}
\end{hintbox}

\label{hint:97}
\begin{hintbox}{97}{Frequenza dei caratteri}
\small\color{black!55}\hyperref[ex:97]{$\leftarrow$ Esercizio}\normalsize\color{black}

Crea una hash table con \fn{:test \#'eql} (i caratteri si
confrontano con \fn{eql}).
Itera sui caratteri con \fn{(coerce stringa 'list)}
o \fn{(loop for c across stringa ...)}.
Per ogni carattere: incrementa il valore nella tabella
(\fn{(incf (gethash c ht 0))}).

\small\color{black!55}\hyperref[sol:97]{$\rightarrow$ Soluzione}\normalsize\color{black}
\end{hintbox}

\label{hint:98}
\begin{hintbox}{98}{Unione di hash table}
\small\color{black!55}\hyperref[ex:98]{$\leftarrow$ Esercizio}\normalsize\color{black}

Crea una nuova hash table. Itera su \fn{ht1} con \fn{maphash}
e copia tutte le coppie. Poi itera su \fn{ht2}: se la chiave
esiste gia', somma i valori; altrimenti inserisci.

\small\color{black!55}\hyperref[sol:98]{$\rightarrow$ Soluzione}\normalsize\color{black}
\end{hintbox}

\label{hint:99}
\begin{hintbox}{99}{Inverti hash table}
\small\color{black!55}\hyperref[ex:99]{$\leftarrow$ Esercizio}\normalsize\color{black}

Crea una nuova hash table. Itera con \fn{maphash} sulla
tabella originale: per ogni coppia \fn{(k v)}, inserisci
\fn{(v k)} nella nuova tabella.

\small\color{black!55}\hyperref[sol:99]{$\rightarrow$ Soluzione}\normalsize\color{black}
\end{hintbox}

\label{hint:102}
\begin{hintbox}{102}{Trasposta di matrice}
\small\color{black!55}\hyperref[ex:102]{$\leftarrow$ Esercizio}\normalsize\color{black}

Crea una nuova matrice $n \times n$. Con due loop annidati
\fn{(loop for i ...)} e \fn{(loop for j ...)}:
\fn{(setf (aref risultato i j) (aref originale j i))}.

\small\color{black!55}\hyperref[sol:102]{$\rightarrow$ Soluzione}\normalsize\color{black}
\end{hintbox}

\label{hint:104}
\begin{hintbox}{104}{Moltiplicazione di matrici (opzionale)}
\small\color{black!55}\hyperref[ex:104]{$\leftarrow$ Esercizio}\normalsize\color{black}

Il prodotto $C = A \times B$ con $C_{ij} = \sum_k A_{ik} \cdot B_{kj}$.
Tre loop annidati: indici $i$, $j$, $k$.
\begin{lstlisting}
(loop for i below n do
  (loop for j below n do
    (setf (aref c i j)
      (loop for k below n
            sum (* (aref a i k)
                   (aref b k j))))))
\end{lstlisting}

\small\color{black!55}\hyperref[sol:104]{$\rightarrow$ Soluzione}\normalsize\color{black}
\end{hintbox}

\label{hint:108}
\begin{hintbox}{108}{BST: inserimento e ricerca}
\small\color{black!55}\hyperref[ex:108]{$\leftarrow$ Esercizio}\normalsize\color{black}

\textbf{Insert:} se l'albero e' \fn{nil}, crea un nodo foglia.
Se il valore e' minore del nodo corrente, ricorri a sinistra;
altrimenti a destra. Costruisci il nuovo albero con \fn{make-nodo}.

\textbf{Find:} se l'albero e' \fn{nil}, non trovato.
Se il valore e' uguale al nodo, trovato.
Se minore, vai a sinistra; se maggiore, vai a destra.

\small\color{black!55}\hyperref[sol:108]{$\rightarrow$ Soluzione}\normalsize\color{black}
\end{hintbox}

\label{hint:109}
\begin{hintbox}{109}{Traversal dell'albero}
\small\color{black!55}\hyperref[ex:109]{$\leftarrow$ Esercizio}\normalsize\color{black}

I tre traversal differiscono solo per l'ordine in cui si
visita la radice rispetto ai sottoalberi:
\begin{itemize}[leftmargin=2em]
  \item \textbf{preorder:} radice, sinistra, destra.
  \item \textbf{inorder:} sinistra, radice, destra.
  \item \textbf{postorder:} sinistra, destra, radice.
\end{itemize}
Usa \fn{append} per concatenare i risultati dei due sottoalberi.

\small\color{black!55}\hyperref[sol:109]{$\rightarrow$ Soluzione}\normalsize\color{black}
\end{hintbox}

\label{hint:112}
\begin{hintbox}{112}{Albero bilanciato?}
\small\color{black!55}\hyperref[ex:112]{$\leftarrow$ Esercizio}\normalsize\color{black}

Scrivi una funzione ausiliaria che restituisce l'altezza
se l'albero e' bilanciato, oppure $-1$ come segnale di
sbilanciamento.
\begin{lstlisting}
; ritorna altezza se bilanciato, -1 altrimenti
(defun check (tree)
  (if (null tree)
      0
      (let ((hl (check (node-left tree)))
            (hr (check (node-right tree))))
        (if (or (= hl -1) (= hr -1)
                (> (abs (- hl hr)) 1))
            -1
            (1+ (max hl hr))))))
\end{lstlisting}

\small\color{black!55}\hyperref[sol:112]{$\rightarrow$ Soluzione}\normalsize\color{black}
\end{hintbox}

\label{hint:114}
\begin{hintbox}{114}{Coda FIFO}
\small\color{black!55}\hyperref[ex:114]{$\leftarrow$ Esercizio}\normalsize\color{black}

Rappresenta la coda come una struttura con due liste:
\fn{testa} (elementi da estrarre) e \fn{coda} (elementi aggiunti).
\begin{itemize}[leftmargin=2em]
  \item \textbf{Enqueue:} aggiungi in \fn{coda} con \fn{cons}.
  \item \textbf{Dequeue:} prendi da \fn{testa}. Se \fn{testa}
        e' vuota, ribalta \fn{coda} in \fn{testa} con \fn{reverse}.
\end{itemize}
Questa rappresentazione da' dequeue ammortizzato $O(1)$.

\small\color{black!55}\hyperref[sol:114]{$\rightarrow$ Soluzione}\normalsize\color{black}
\end{hintbox}

\label{hint:115}
\begin{hintbox}{115}{Valida parentesi}
\small\color{black!55}\hyperref[ex:115]{$\leftarrow$ Esercizio}\normalsize\color{black}

Itera sui caratteri della stringa. Per ogni carattere:
\begin{itemize}[leftmargin=2em]
  \item parentesi aperta \fn{(}, \fn{[}, \fn{\{}: push sulla pila.
  \item parentesi chiusa \fn{)}, \fn{]}, \fn{\}}: se la pila
        e' vuota o il top non corrisponde, restituisci \fn{NIL}.
        Altrimenti pop.
\end{itemize}
Alla fine, la stringa e' valida se e solo se la pila e' vuota.
Usa una alist per mappare ogni chiusa alla corrispondente aperta.

\small\color{black!55}\hyperref[sol:115]{$\rightarrow$ Soluzione}\normalsize\color{black}
\end{hintbox}

\label{hint:118}
\begin{hintbox}{118}{DFS su grafo}
\small\color{black!55}\hyperref[ex:118]{$\leftarrow$ Esercizio}\normalsize\color{black}

Usa una lista \fn{visitati} per tenere traccia dei nodi gia' visti.
Schema ricorsivo:
\begin{enumerate}[leftmargin=2em]
  \item Se il nodo e' gia' visitato, restituisci \fn{NIL}.
  \item Aggiungi il nodo a \fn{visitati}.
  \item Per ogni vicino non ancora visitato, ricorri.
  \item Accumula l'ordine di visita.
\end{enumerate}
Passa \fn{visitati} come parametro o usala come variabile
catturata da una closure/\fn{labels}.

\small\color{black!55}\hyperref[sol:118]{$\rightarrow$ Soluzione}\normalsize\color{black}
\end{hintbox}

\label{hint:119}
\begin{hintbox}{119}{BFS su grafo}
\small\color{black!55}\hyperref[ex:119]{$\leftarrow$ Esercizio}\normalsize\color{black}

Usa una coda (o semplicemente una lista come coda FIFO).
Schema:
\begin{enumerate}[leftmargin=2em]
  \item Inserisci il nodo di partenza nella coda.
  \item Finche' la coda non e' vuota:
        estrai il primo nodo; se non ancora visitato,
        aggiungilo ai visitati e inserisci i suoi vicini in coda.
\end{enumerate}
Con una lista, usa \fn{append} per aggiungere in coda e
\fn{car}/\fn{cdr} per estrarre dalla testa.

\small\color{black!55}\hyperref[sol:119]{$\rightarrow$ Soluzione}\normalsize\color{black}
\end{hintbox}

\label{hint:120}
\begin{hintbox}{120}{Conta parole}
\small\color{black!55}\hyperref[ex:120]{$\leftarrow$ Esercizio}\normalsize\color{black}

Per spezzare la stringa in parole, implementa un tokenizer
semplice: scansiona carattere per carattere, accumulando
caratteri non-spazio in una parola corrente.
Quando trovi uno spazio (o la fine), se la parola corrente
non e' vuota, aggiungila alla lista.
Poi usa una hash table con \fn{:test \#'equal} (le stringhe
si confrontano con \fn{equal}) per contare le occorrenze.

\small\color{black!55}\hyperref[sol:120]{$\rightarrow$ Soluzione}\normalsize\color{black}
\end{hintbox}

% =============================================================
\hintsection{Macro}
% =============================================================

\label{hint:121}
\begin{hintbox}{121}{Macro quando}
\small\color{black!55}\hyperref[ex:121]{$\leftarrow$ Esercizio}\normalsize\color{black}

Lo schema base di una macro con body:
\begin{lstlisting}
(defmacro quando (test &body corpo)
  `(if ,test
       (progn ,@corpo)
       nil))
\end{lstlisting}
Il backquote costruisce la lista quasi-letteralmente;
\fn{,test} inserisce il valore di \fn{test};
\fn{,@corpo} inserisce gli elementi di \fn{corpo}
al livello corrente (splice).

Verifica con: \fn{(macroexpand-1 '(quando (> x 0) (print x)))}.

\small\color{black!55}\hyperref[sol:121]{$\rightarrow$ Soluzione}\normalsize\color{black}
\end{hintbox}

\label{hint:122}
\begin{hintbox}{122}{Macro while}
\small\color{black!55}\hyperref[ex:122]{$\leftarrow$ Esercizio}\normalsize\color{black}

Un'implementazione semplice usa \fn{loop} nell'espansione:
\begin{lstlisting}
(defmacro mio-while (test &body corpo)
  `(loop while ,test do ,@corpo))
\end{lstlisting}
Oppure, per capire meglio il meccanismo di basso livello,
usa \fn{tagbody}/\fn{go}:
\begin{lstlisting}
(defmacro mio-while (test &body corpo)
  (let ((inizio (gensym)) (fine (gensym)))
    `(tagbody
       ,inizio
       (unless ,test (go ,fine))
       ,@corpo
       (go ,inizio)
       ,fine)))
\end{lstlisting}
\fn{gensym} genera simboli unici per evitare conflitti di nomi.

\small\color{black!55}\hyperref[sol:122]{$\rightarrow$ Soluzione}\normalsize\color{black}
\end{hintbox}

\label{hint:123}
\begin{hintbox}{123}{Macro for}
\small\color{black!55}\hyperref[ex:123]{$\leftarrow$ Esercizio}\normalsize\color{black}

Usa \fn{gensym} per la variabile di ciclo e per evitare che
valutazioni multiple di \fn{da}/\fn{a} causino effetti collaterali.
\begin{lstlisting}
(defmacro mio-for ((var da a) &body corpo)
  (let ((start (gensym)) (end (gensym)))
    `(let ((,start ,da) (,end ,a))
       (loop for ,var from ,start to ,end
             do ,@corpo))))
\end{lstlisting}
Nota: qui \fn{var} non e' avvolto in \fn{gensym} perche'
e' il nome che l'utente vuole usare nel corpo.

\small\color{black!55}\hyperref[sol:123]{$\rightarrow$ Soluzione}\normalsize\color{black}
\end{hintbox}

% =============================================================
\hintsection{Backtracking e algoritmi}
% =============================================================

\label{hint:124}
\begin{hintbox}{124}{Subsets}
\small\color{black!55}\hyperref[ex:124]{$\leftarrow$ Esercizio}\normalsize\color{black}

Ogni elemento ha due scelte: incluso o non incluso.
Schema ricorsivo:
\begin{itemize}[leftmargin=2em]
  \item Caso base: lista vuota $\to$ \fn{'(())}.
  \item Caso ricorsivo: calcola i sottoinsiemi del \fn{cdr}.
        I sottoinsiemi totali sono quelli senza \fn{(car lst)},
        piu' quelli con \fn{(car lst)} aggiunto in testa.
\end{itemize}

\small\color{black!55}\hyperref[sol:124]{$\rightarrow$ Soluzione}\normalsize\color{black}
\end{hintbox}

\label{hint:125}
\begin{hintbox}{125}{Permutazioni}
\small\color{black!55}\hyperref[ex:125]{$\leftarrow$ Esercizio}\normalsize\color{black}

Per generare le permutazioni di \fn{lst}:
\begin{enumerate}[leftmargin=2em]
  \item Per ogni elemento \fn{x} nella lista:
  \item Rimuovi \fn{x} dalla lista (ottenendo \fn{rest}).
  \item Genera le permutazioni di \fn{rest}.
  \item Aggiungi \fn{x} in testa a ogni permutazione.
\end{enumerate}
Usa una funzione ausiliaria che rimuove un elemento per posizione
(non per valore, per gestire duplicati correttamente).

\small\color{black!55}\hyperref[sol:125]{$\rightarrow$ Soluzione}\normalsize\color{black}
\end{hintbox}

\label{hint:126}
\begin{hintbox}{126}{Combinazioni}
\small\color{black!55}\hyperref[ex:126]{$\leftarrow$ Esercizio}\normalsize\color{black}

Per scegliere \fn{k} elementi da \fn{lst}:
\begin{itemize}[leftmargin=2em]
  \item Se \fn{k = 0}: restituisci \fn{'(())}.
  \item Se la lista e' vuota: \fn{NIL}.
  \item Altrimenti: unisci le combinazioni che \textit{includono}
        \fn{(car lst)} (combinazioni di \fn{k-1} dal resto)
        con quelle che \textit{non lo includono} (combinazioni
        di \fn{k} dal resto).
\end{itemize}

\small\color{black!55}\hyperref[sol:126]{$\rightarrow$ Soluzione}\normalsize\color{black}
\end{hintbox}

\label{hint:127}
\begin{hintbox}{127}{Somma a target}
\small\color{black!55}\hyperref[ex:127]{$\leftarrow$ Esercizio}\normalsize\color{black}

Adatta il pattern dei sottoinsiemi: a ogni passo, scegli
di includere o meno il numero corrente.
Se il \fn{target} diventa 0, hai trovato una soluzione.
Se il \fn{target} diventa negativo o la lista e' vuota
senza aver raggiunto 0, il ramo fallisce.

\small\color{black!55}\hyperref[sol:127]{$\rightarrow$ Soluzione}\normalsize\color{black}
\end{hintbox}

\label{hint:128}
\begin{hintbox}{128}{Parentesi bilanciate generate}
\small\color{black!55}\hyperref[ex:128]{$\leftarrow$ Esercizio}\normalsize\color{black}

Costruisci la stringa con due contatori: \fn{aperte} e \fn{chiuse}.
A ogni passo puoi aggiungere \fn{(} se \fn{aperte < n};
puoi aggiungere \fn{)} se \fn{chiuse < aperte}.
Quando \fn{aperte = chiuse = n}, hai una stringa valida.

\small\color{black!55}\hyperref[sol:128]{$\rightarrow$ Soluzione}\normalsize\color{black}
\end{hintbox}

\label{hint:129}
\begin{hintbox}{129}{Merge sort}
\small\color{black!55}\hyperref[ex:129]{$\leftarrow$ Esercizio}\normalsize\color{black}

\textbf{Merge-sorted:} due puntatori alle teste delle liste.
Prendi il minore dei due \fn{car}, avanza il puntatore corrispondente.
Caso base: quando una lista e' vuota, aggiungi il resto dell'altra.

\textbf{Merge-sort:} dividi la lista a meta' (usa \fn{length} e
\fn{nthcdr} per trovare il punto di divisione),
ordina le due meta' ricorsivamente, poi unisci con \fn{merge-sorted}.

\small\color{black!55}\hyperref[sol:129]{$\rightarrow$ Soluzione}\normalsize\color{black}
\end{hintbox}

\label{hint:130}
\begin{hintbox}{130}{Quicksort}
\small\color{black!55}\hyperref[ex:130]{$\leftarrow$ Esercizio}\normalsize\color{black}

Scegli il primo elemento come pivot.
Con \fn{remove-if} (o la tua \fn{filter-list}), dividi
il resto in elementi minori e maggiori-o-uguali al pivot.
Ordina ricorsivamente le due parti e concatena:
\fn{minori ++ [pivot] ++ maggiori}.

\small\color{black!55}\hyperref[sol:130]{$\rightarrow$ Soluzione}\normalsize\color{black}
\end{hintbox}

\label{hint:131}
\begin{hintbox}{131}{Torri di Hanoi}
\small\color{black!55}\hyperref[ex:131]{$\leftarrow$ Esercizio}\normalsize\color{black}

La soluzione e' elegantemente ricorsiva:
\begin{enumerate}[leftmargin=2em]
  \item Sposta $n-1$ dischi da \fn{da} a \fn{tramite} (usando \fn{a}).
  \item Sposta il disco $n$ da \fn{da} a \fn{a}.
  \item Sposta $n-1$ dischi da \fn{tramite} a \fn{a} (usando \fn{da}).
\end{enumerate}
Caso base: \fn{n = 0}, non fare nulla.
Il numero di mosse e' $2^n - 1$.

\small\color{black!55}\hyperref[sol:131]{$\rightarrow$ Soluzione}\normalsize\color{black}
\end{hintbox}

\label{hint:132}
\begin{hintbox}{132}{Problema dello zaino}
\small\color{black!55}\hyperref[ex:132]{$\leftarrow$ Esercizio}\normalsize\color{black}

Schema di ricerca esaustiva:
Per ogni oggetto, hai due scelte: prenderlo o lasciarlo.
\begin{itemize}[leftmargin=2em]
  \item Prenderlo: sottrai il peso dalla capacita', aggiungi
        il valore, ricorri sul resto degli oggetti.
  \item Non prenderlo: ricorri sul resto con capacita' invariata.
\end{itemize}
Restituisci il massimo tra i due rami.
Attenzione: prendi l'oggetto solo se il suo peso non supera
la capacita' residua.

\small\color{black!55}\hyperref[sol:132]{$\rightarrow$ Soluzione}\normalsize\color{black}
\end{hintbox}

\label{hint:133}
\begin{hintbox}{133}{Labirinto}
\small\color{black!55}\hyperref[ex:133]{$\leftarrow$ Esercizio}\normalsize\color{black}

Schema DFS con backtracking:
\begin{enumerate}[leftmargin=2em]
  \item Se la cella corrente e' il goal, restituisci il percorso.
  \item Marca la cella come visitata.
  \item Per ogni direzione (su, giu', sinistra, destra):
        se la cella adiacente e' valida, libera e non visitata,
        ricorri. Se restituisce un percorso, propagalo.
  \item Se nessuna direzione funziona, demarks e restituisci \fn{NIL}.
\end{enumerate}
Usa un array separato per i visitati, oppure marca nel labirinto stesso
(ricordando di ripristinare al backtrack).

\small\color{black!55}\hyperref[sol:133]{$\rightarrow$ Soluzione}\normalsize\color{black}
\end{hintbox}

\label{hint:134}
\begin{hintbox}{134}{N-regine}
\small\color{black!55}\hyperref[ex:134]{$\leftarrow$ Esercizio}\normalsize\color{black}

\hintschema{Rappresentazione:} usa un vettore di \fn{n} elementi dove
\fn{(aref board col)} e' la riga in cui si trova la regina
nella colonna \fn{col}. Stai posizionando una regina per colonna.

\hintschema{Pruning:} prima di posizionare la regina nella colonna \fn{col}
alla riga \fn{row}, verifica che nessuna delle regine gia'
posizionate attacchi questa cella:
\begin{itemize}[leftmargin=2em]
  \item stessa riga: \fn{(= (aref board prev-col) row)};
  \item stessa diagonale: \fn{(= (abs (- row (aref board prev-col))) (- col prev-col))}.
\end{itemize}

\textbf{Struttura:} per ogni colonna, prova tutte le righe 0..(n-1).
Se valida, assegna e ricorri sulla colonna successiva.
Se \fn{col = n}, hai trovato una soluzione.

\small\color{black!55}\hyperref[sol:134]{$\rightarrow$ Soluzione}\normalsize\color{black}
\end{hintbox}

\label{hint:135}
\begin{hintbox}{135}{Sudoku}
\small\color{black!55}\hyperref[ex:135]{$\leftarrow$ Esercizio}\normalsize\color{black}

\begin{enumerate}[leftmargin=2em]
  \item Trova la prima cella vuota (valore 0).
  \item Se non ce ne sono, il puzzle e' risolto.
  \item Per ogni valore da 1 a 9:
        verifica che non compaia nella stessa riga, colonna,
        e blocco $3 \times 3$.
        Se valido, assegna il valore e ricorri.
        Se la ricorsione fallisce, rimetti 0 (backtrack).
\end{enumerate}

Per ottimizzare, scegli la cella con meno valori possibili
(``minimum remaining values heuristic'').

\small\color{black!55}\hyperref[sol:135]{$\rightarrow$ Soluzione}\normalsize\color{black}
\end{hintbox}

\label{hint:136}
\begin{hintbox}{136}{Cammino del cavallo}
\small\color{black!55}\hyperref[ex:136]{$\leftarrow$ Esercizio}\normalsize\color{black}

Le 8 mosse possibili del cavallo in termini di $(dx, dy)$:
$(\pm 1, \pm 2)$ e $(\pm 2, \pm 1)$.

\textbf{Backtracking puro:} prova ogni mossa valida (dentro la
scacchiera e non gia' visitata); se porta a una soluzione,
propagala; altrimenti annulla e prova la successiva.

\textbf{Euristica di Warnsdorff:} ordina le mosse in base al numero
di uscite disponibili dalla cella successiva (meno uscite = prima).
Riduce drasticamente il backtracking.

\small\color{black!55}\hyperref[sol:136]{$\rightarrow$ Soluzione}\normalsize\color{black}
\end{hintbox}

% =============================================================
\hintsection{Stringhe e I/O}
% =============================================================

\label{hint:139}
\begin{hintbox}{139}{Palindromo robusto}
\small\color{black!55}\hyperref[ex:139]{$\leftarrow$ Esercizio}\normalsize\color{black}

\begin{enumerate}[leftmargin=2em]
  \item Converti in minuscolo con \fn{string-downcase}.
  \item Filtra i soli caratteri alfabetici con \fn{remove-if-not}
        e \fn{alpha-char-p}.
  \item Confronta la stringa filtrata con la sua inversa.
\end{enumerate}

\small\color{black!55}\hyperref[sol:139]{$\rightarrow$ Soluzione}\normalsize\color{black}
\end{hintbox}

\label{hint:140}
\begin{hintbox}{140}{Anagrammi}
\small\color{black!55}\hyperref[ex:140]{$\leftarrow$ Esercizio}\normalsize\color{black}

Due stringhe sono anagrammi se, ordinate lettera per lettera,
sono uguali. Converti entrambe in minuscolo, poi in lista
di caratteri con \fn{coerce}, ordina con \fn{sort}, e confronta.

\small\color{black!55}\hyperref[sol:140]{$\rightarrow$ Soluzione}\normalsize\color{black}
\end{hintbox}

% =============================================================
\hintsection{Parsing e interprete}
% =============================================================

\label{hint:142}
\begin{hintbox}{142}{Calcolatrice RPN}
\small\color{black!55}\hyperref[ex:142]{$\leftarrow$ Esercizio}\normalsize\color{black}

Itera sugli elementi dell'espressione:
\begin{itemize}[leftmargin=2em]
  \item Se e' un numero: push sulla pila.
  \item Se e' un operatore: pop due valori, applica
        l'operatore, push il risultato.
\end{itemize}
Alla fine, la pila contiene un solo elemento: il risultato.
Usa \fn{numberp} per distinguere numeri da operatori;
usa una alist o \fn{case} per mappare simbolo $\to$ funzione.

\small\color{black!55}\hyperref[sol:142]{$\rightarrow$ Soluzione}\normalsize\color{black}
\end{hintbox}

\label{hint:143}
\begin{hintbox}{143}{Valutatore di espressioni}
\small\color{black!55}\hyperref[ex:143]{$\leftarrow$ Esercizio}\normalsize\color{black}

La struttura dell'espressione rispecchia la struttura del codice:
\begin{itemize}[leftmargin=2em]
  \item Se e' un numero: restituiscilo direttamente.
  \item Se e' una lista: il \fn{car} e' l'operatore,
        il \fn{cadr} e il \fn{caddr} sono i due operandi.
        Valutali ricorsivamente, poi applica l'operatore.
\end{itemize}
Usa \fn{funcall} con una funzione ricavata dal simbolo
(es. \fn{(symbol-function operator)}) oppure un \fn{case}.

\small\color{black!55}\hyperref[sol:143]{$\rightarrow$ Soluzione}\normalsize\color{black}
\end{hintbox}

\label{hint:144}
\begin{hintbox}{144}{Mini-interprete con variabili}
\small\color{black!55}\hyperref[ex:144]{$\leftarrow$ Esercizio}\normalsize\color{black}

L'ambiente e' una alist \fn{((nome . valore) ...)}.

Estendi il valutatore dell'es.~143:
\begin{itemize}[leftmargin=2em]
  \item Se l'espressione e' un simbolo: cercalo nell'ambiente
        con \fn{assoc}.
  \item Se e' \fn{(set var val)}: valuta \fn{val}, aggiorna
        l'ambiente (aggiungi o modifica la coppia).
  \item Operatori aritmetici: come prima.
\end{itemize}
Gestisci l'ambiente come parametro (passato per riferimento
o restituito insieme al valore).

\small\color{black!55}\hyperref[sol:144]{$\rightarrow$ Soluzione}\normalsize\color{black}
\end{hintbox}

% =============================================================
\hintsection{Progetti}
% =============================================================

\label{hint:146}
\begin{hintbox}{146}{Tris}
\small\color{black!55}\hyperref[ex:146]{$\leftarrow$ Esercizio}\normalsize\color{black}

Struttura consigliata:
\begin{itemize}[leftmargin=2em]
  \item \fn{crea-griglia}: array $3 \times 3$ di \fn{NIL}.
  \item \fn{fai-mossa griglia riga col giocatore}: verifica
        che la cella sia libera e assegna.
  \item \fn{vinto-p griglia giocatore}: controlla le 8
        combinazioni vincenti (3 righe, 3 colonne, 2 diagonali).
  \item \fn{pareggio-p griglia}: nessuna cella vuota e nessuno ha vinto.
  \item \fn{gioca}: loop che alterna i turni fino a vittoria o pareggio.
\end{itemize}

\small\color{black!55}\hyperref[sol:146]{$\rightarrow$ Soluzione}\normalsize\color{black}
\end{hintbox}

\label{hint:147}
\begin{hintbox}{147}{Gioco della vita}
\small\color{black!55}\hyperref[ex:147]{$\leftarrow$ Esercizio}\normalsize\color{black}

Per ogni cella $(r, c)$, conta i vicini vivi (le 8 celle adiacenti).
Le regole:
\begin{itemize}[leftmargin=2em]
  \item Viva con 2 o 3 vicini $\to$ sopravvive.
  \item Morta con esattamente 3 vicini $\to$ nasce.
  \item Altrimenti $\to$ muore/resta morta.
\end{itemize}
Calcola la nuova generazione in un array separato (non
modificare quello corrente mentre lo leggi).

\small\color{black!55}\hyperref[sol:147]{$\rightarrow$ Soluzione}\normalsize\color{black}
\end{hintbox}

\label{hint:148}
\begin{hintbox}{148}{Mini sistema gestionale}
\small\color{black!55}\hyperref[ex:148]{$\leftarrow$ Esercizio}\normalsize\color{black}

Struttura minima consigliata:

\textbf{Model:} definisci la struttura con \fn{defstruct}.

\textbf{Archivio:} una hash table globale \fn{*archivio*}
con chiave = identificatore univoco.

\textbf{Salvataggio:} con \fn{with-open-file} scrivi ogni
record come S-expression con \fn{print} (leggibile con \fn{read}).
\begin{lstlisting}
(with-open-file (s file :direction :output
                        :if-exists :supersede)
  (maphash (lambda (k v)
             (print (list k v) s))
           *archivio*))
\end{lstlisting}

\textbf{Caricamento:} leggi le S-expression con \fn{read}
in un loop finche' non si raggiunge \fn{EOF}.

\small\color{black!55}\hyperref[sol:148]{$\rightarrow$ Soluzione}\normalsize\color{black}
\end{hintbox}

